\documentclass[10pt,a4paper]{amsart}
\usepackage[margin=19mm]{geometry}
\usepackage{amsmath,amssymb,mathtools}
\usepackage[scr=rsfs,cal=euler]{mathalfa}
\usepackage{xfrac}
\usepackage[unicode,hidelinks,bookmarks=false]{hyperref}
\newcommand{\ie}{i.e.\ }
\newcommand{\cf}{cf.\ }
\newcommand{\resp}{respectively\ }
\newcommand{\Subsection}{\subsection}
\providecommand{\nobreakdash}{\nobreak\hbox{-}}
\setlength{\emergencystretch}{2em}
\multlinegap0pt
\usepackage[shortlabels]{enumitem}
\usepackage{multirow}
\usepackage{amssymb}
\usepackage{mathtools}
\usepackage{xcolor}
\usepackage{booktabs}
\usepackage{algorithm}\usepackage{algpseudocode}
\usepackage{tikz}
\newtheorem{theorem}{Theorem}
\title{Beyond Diagonal Noise: A Better Predator-Prey Modeling Framework with Cross-Covariance}
\author{Jiguang Yu, Louis Shuo Wang}
\date{Source: arXiv:2602.22489v1 (CC BY 4.0)}
\begin{document}
\maketitle

\noindent\emph{Source and licence.} Beyond Diagonal Noise: A Better Predator-Prey Modeling Framework with Cross-Covariance. Version arXiv:2602.22489v1 (CC BY 4.0), distributed by arXiv under the Creative Commons Attribution 4.0 International licence, http://creativecommons.org/licenses/by/4.0/, as recorded in the arXiv licence metadata for this exact version and shown on the abstract page https://arxiv.org/abs/2602.22489v1. Full licence notice: \texttt{licenses/arxiv-2602.22489-CC-BY-4.0.txt}.\par\smallskip

\noindent\textit{Editorial scope.} Counted unit: the article's theorem thm:maximal\_no\_interior\_explosion\_paperA (source0.tex lines 971--1002). The article's complete original proof of it is delivered with it (source0.tex lines 1636--1716, one \texttt{proof} environment of 81 lines, which the article places away from the statement). No mathematics is written, repaired or re-derived: the statement and the proof are the article's own lines, in the article's own notation, and no retained line is substituted or re-flowed.  Retained uncounted as the source-local context the proof needs: lines 904--908; lines 1459--1462; lines 1468--1471.  Nothing was omitted from any retained span, so the package carries no omission record.  No retained line contains a citation, so the wrapper carries no bibliography and no bibliography entry.  No retained line contains a figure environment, an image-inclusion command or drawing code, so no image is delivered and the package declares no asset dependency.  Editorial formatting only, in addition: an AMS \texttt{amsart} preamble that restates the article's own declarations and macros used by the retained lines, namely the environments \texttt{theorem} on separate counters, together with the standard packages those lines need.  The retained environments print in this document as Theorem 1 in order of appearance, whereas the article prints the same items with the numbering of its own full document.  The complete native source of the article as distributed in the arXiv e-print for this exact version is bundled unchanged as \texttt{sources/195336/source0.tex}.
\begin{equation}
\label{eq:open_domain_SDE_paperA}
dx_t=b(x_t)\,dt+\Sigma(x_t)\,dW_t,
\qquad x_0\in D.
\end{equation}

\begin{equation}
\label{eq:appC_exhaustion}
K_n\subset \operatorname{int}(K_{n+1}),\qquad K_n\Subset D,\qquad \bigcup_{n=1}^\infty K_n=D.
\end{equation}

\begin{equation}
\label{eq:appC_tau_n}
\tau_n:=\inf\{t\ge 0:\ x_t\notin K_n\}.
\end{equation}

\begin{theorem}[Maximal strong solution and no interior explosion before boundary hitting]
\label{thm:maximal_no_interior_explosion_paperA}
Assume:
\begin{enumerate}[label=(A\arabic*)]
\item \textbf{Local regularity on compacts.} For every precompact set \(K\Subset D\), the coefficients \(b\) and \(\Sigma\) are Lipschitz and bounded on \(K\).
\item \textbf{Lyapunov control at infinity.} There exist \(V_\infty\in C^2(D;[0,\infty))\) and a constant \(C>0\) such that
\begin{equation}
\label{eq:Lyap_infty_cond_paperA}
\mathcal L V_\infty(x)\le C\bigl(1+V_\infty(x)\bigr),\qquad x\in D,
\end{equation}
where
\[
\mathcal L V(x)=\nabla V(x)\cdot b(x)+\frac12\mathrm{tr}\!\bigl(a(x)\nabla^2V(x)\bigr),
\qquad a(x)=\Sigma(x)\Sigma(x)^\top,
\]
and \(V_\infty(x)\to\infty\) whenever \(\|x\|\to\infty\) with \(x\in D\).
\end{enumerate}
Then for every \(x_0\in D\), the SDE~\eqref{eq:open_domain_SDE_paperA} admits a unique maximal strong solution
\[
(x_t)_{0\le t<\tau_*}
\]
with lifetime \(\tau_*\in(0,\infty]\), and
\begin{equation}
\label{eq:no_interior_blowup_paperA}
\mathbb P\bigl(\tau_*<\infty,\ \tau_*<\tau_{\partial D}\bigr)=0.
\end{equation}
Equivalently, loss of well-posedness cannot occur in the interior of \(D\); if \(\tau_*<\infty\), the only possible obstruction is boundary hitting (or boundary approach), i.e.
\begin{equation}
\label{eq:tau_star_equals_boundary_paperA}
\tau_*=\tau_{\partial D}\qquad\text{a.s.}
\end{equation}
\end{theorem}

\begin{proof}
Let \((K_n)\) and \((\tau_n)\) be the exhaustion and exit times defined in \eqref{eq:appC_exhaustion}--\eqref{eq:appC_tau_n}. By Theorem~\ref{thm:maximal_no_interior_explosion_paperA}, a unique maximal strong solution exists up to \(\tau_*\), and any finite lifetime obstruction can only occur through boundary hitting/approach.

\medskip
\noindent\textbf{Step 1: Localized It\^o estimate for the barrier Lyapunov function.}
Fix \(n\ge1\) and \(t\ge0\). Since \(x_{s\wedge \tau_n}\in K_n\) for all \(s\), It\^o's formula for \(V_b(x_{t\wedge \tau_n})\) yields
\begin{align}
V_b(x_{t\wedge \tau_n})
&=
V_b(x_0)+\int_0^{t\wedge \tau_n}\mathcal LV_b(x_s)\,ds + M_{t\wedge \tau_n}^{(b)} \notag\\
&\le
V_b(x_0)+C_b\int_0^{t\wedge \tau_n}\bigl(1+V_b(x_s)\bigr)\,ds + M_{t\wedge \tau_n}^{(b)},
\label{eq:appC_barrier_ito}
\end{align}
where \(M_{t\wedge \tau_n}^{(b)}\) is a true martingale.

Taking expectations,
\begin{equation}
\label{eq:appC_barrier_expect_pre_gronwall}
\mathbb E[V_b(x_{t\wedge \tau_n})]
\le
V_b(x_0)+C_b t + C_b\int_0^t \mathbb E[V_b(x_{s\wedge \tau_n})]\,ds.
\end{equation}
Applying Gr\"onwall,
\begin{equation}
\label{eq:appC_barrier_expect_bound}
\mathbb E[V_b(x_{t\wedge \tau_n})]
\le e^{C_b t}\bigl(V_b(x_0)+C_b t\bigr),\qquad \forall n\ge1,\ \forall t\ge0.
\end{equation}
Hence
\begin{equation}
\label{eq:appC_barrier_uniform_bound}
\sup_{n\ge1}\mathbb E[V_b(x_{t\wedge \tau_n})]<\infty\qquad \text{for each fixed }t\ge0.
\end{equation}

\medskip
\noindent\textbf{Step 2: Contradiction argument if boundary is hit with positive probability.}
Assume for contradiction that \(\mathbb P(\tau_{\partial D}<\infty)>0\). Then there exists \(T>0\) such that
\[
\mathbb P(\tau_{\partial D}\le T)>0.
\]
On the event \(\{\tau_{\partial D}\le T\}\), continuity of sample paths implies \(x_t\to x_{\tau_{\partial D}}\in\partial D\) as \(t\uparrow \tau_{\partial D}\) from below (or equivalently \(x_t\to \partial D\) within \(D\)). By the barrier blow-up property \textbf{(B1)},
\[
V_b(x_t)\to\infty\qquad \text{as }t\uparrow \tau_{\partial D}\ \text{on }\{\tau_{\partial D}\le T\}.
\]
Because \(\tau_n\uparrow \tau_*\) and Theorem~\ref{thm:maximal_no_interior_explosion_paperA} implies \(\tau_*\ge \tau_{\partial D}\) a.s., the stopping sequence \(T\wedge \tau_n\) approaches \(T\wedge \tau_{\partial D}\) from below on \(\{\tau_{\partial D}\le T\}\). In particular,
\begin{equation}
\label{eq:appC_barrier_divergence}
V_b\bigl(x_{T\wedge \tau_n}\bigr)\to\infty
\qquad\text{on }\{\tau_{\partial D}\le T\}.
\end{equation}
Applying Fatou's lemma to the nonnegative random variables \(V_b(x_{T\wedge \tau_n})\mathbf 1_{\{\tau_{\partial D}\le T\}}\), we obtain
\begin{align}
\infty
&=
\mathbb E\!\left[\liminf_{n\to\infty}V_b(x_{T\wedge \tau_n})\mathbf 1_{\{\tau_{\partial D}\le T\}}\right] \notag\\
&\le
\liminf_{n\to\infty}\mathbb E\!\left[V_b(x_{T\wedge \tau_n})\mathbf 1_{\{\tau_{\partial D}\le T\}}\right]
\le
\liminf_{n\to\infty}\mathbb E[V_b(x_{T\wedge \tau_n})],
\label{eq:appC_fatou_contradiction}
\end{align}
contradicting the uniform bound \eqref{eq:appC_barrier_uniform_bound} at time \(T\).

Therefore \(\mathbb P(\tau_{\partial D}\le T)=0\) for every \(T>0\), hence
\[
\mathbb P(\tau_{\partial D}<\infty)=0.
\]

\medskip
\noindent\textbf{Step 3: Globality and positivity invariance.}
Since boundary hitting is a.s.\ impossible and Theorem~\ref{thm:maximal_no_interior_explosion_paperA} excludes interior breakdown before boundary hitting, the maximal solution extends for all times:
\[
\tau_*=\infty\qquad\text{a.s.}
\]
Moreover,
\[
x_t\in D\qquad \text{for all }t\ge0\quad\text{a.s.}
\]
This proves positivity invariance.
\end{proof}

\end{document}
